% !PRINTED 386
\weberpage{386}

%<<CONT
then, since $-1:\omega'$ and $-1:\omega_1'$ are also reduced, $\alpha$ is the greatest integer contained in $-1:\omega_1'$, and $\omega'$ is completely determined by $\omega_1'$ according to (9), and hence also $\omega$ by $\omega_1$.

From this it follows:

\gesperrt{3. If a reduced number $\omega$ is developed in a continued fraction, then all the closing numbers $\omega_n$ are again reduced numbers, and by any one of these closing numbers $\omega_n$ both the following $\omega_{n+1}$ and the preceding $\omega_{n-1}$ are completely determined.}

\begin{center}
\S.\,125.

\gesperrt{Development of real quadratic irrational numbers}\\
\gesperrt{in continued fractions.}
\end{center}

\gesperrt{4. For a given discriminant there is only a finite number of reduced numbers.}

The correctness of this theorem is easily seen if, according to \S.\,122, one puts the irrational number $\omega$ in the form:

\weq{1}{\omega = \frac{b+\sqrt{D}}{2c} = \frac{-2a}{b-\sqrt{D}}, \quad D = b^2 + 4ac.}

The conjugate number is

\weq{2}{\omega' = \frac{b-\sqrt{D}}{2c} = \frac{-2a}{b+\sqrt{D}}.}

If $\omega$ is to be reduced, then

\weq{3}{0 < \frac{\sqrt{D}-b}{2c} < 1 < \frac{\sqrt{D}+b}{2c}.}

Hence $\sqrt{D}$ must have the same sign as $c$; if we take it to be positive, then

\weq{4}{0 < \sqrt{D} - b < 2c < \sqrt{D} + b.}

A second form of these conditions we obtain from the second representation

\weq{5}{0 < \sqrt{D} - b < 2a < \sqrt{D} + b.}

Hence $a$, $b$, $c$ are positive and $b$ is less than $\sqrt{D}$.

If $\lambda$ denotes the greatest integer contained in $\sqrt{D}$, so that
\[
\lambda < \sqrt{D} < \lambda + 1,
\]
%CONT>>
